\documentclass[10pt,a4paper]{amsart}
\usepackage[margin=19mm]{geometry}
\usepackage{amsmath,amssymb,mathtools}
\usepackage[scr=rsfs,cal=euler]{mathalfa}
\usepackage{xfrac}
\usepackage[unicode,hidelinks,bookmarks=false]{hyperref}
\newcommand{\ie}{i.e.\ }
\newcommand{\cf}{cf.\ }
\newcommand{\resp}{respectively\ }
\newcommand{\Subsection}{\subsection}
\providecommand{\nobreakdash}{\nobreak\hbox{-}}
\setlength{\emergencystretch}{2em}
\multlinegap0pt
\usepackage{bm}
\usepackage{bbm}
\usepackage{dsfont}
\usepackage{xspace}
\usepackage{cancel}
\usepackage{mathtools}
\usepackage{amsthm}
\makeatletter
\@ifundefined{Th}{\newtheorem{Th}{Th}}{}
\@ifundefined{Claim}{\newtheorem{Claim}[Th]{Claim}}{}
\@ifundefined{Cor}{\newtheorem{Cor}[Th]{Corollary}}{}
\@ifundefined{Lemma}{\newtheorem{Lemma}[Th]{Lemma}}{}
\@ifundefined{Prop}{\newtheorem{Prop}[Th]{Proposition}}{}
\makeatother
\title[Area bounds for constant mean curvature surfaces in hyperbol]{Area bounds for constant mean curvature surfaces in hyperbolic 3-manifolds}
\author{Ruojing Jiang}
\date{Source: arXiv:2509.09881v1 (CC BY 4.0)}
\begin{document}
\maketitle

\noindent\emph{Source and licence.} Area bounds for constant mean curvature surfaces in hyperbolic 3-manifolds. Version arXiv:2509.09881v1 (CC BY 4.0), distributed under the Creative Commons Attribution 4.0 International licence, as recorded in the arXiv licence metadata for this exact version and shown on the abstract page https://arxiv.org/abs/2509.09881v1. Full rights notice: licenses/arxiv-2509.09881-CC-BY-4.0.txt.\par\smallskip

\noindent\textit{Editorial scope.} Counted unit: the Prop whose source label is \texttt{prop\_extrinsic curv} in the article (source0.tex lines 311--316, source label \texttt{prop\_extrinsic curv}), delivered together with the article's own complete original proof of it (source0.tex lines 319--363, one \texttt{proof} environment of 45 lines). No mathematics is written, repaired or re-derived: statement and proof are the article's own text line for line. Required context retained uncounted: lemma\_weak extrinsic (lines 229--233); cor\_extrinsic radius (lines 264--268). Every definition, display and label of those lines is retained unchanged. Omitted from this package and not redistributed: 1--228, 234--263, 269--310, 317--318, 364--726. Nothing inside a retained excerpt is omitted or altered: every retained body is delivered exactly as the article's lines are written, so no omission receipt of any kind accompanies this package. Citations: the retained excerpts carry no citation command at all, so no bibliography entry of the article is transcribed and no reference is redistributed. Reference adaptation: none was needed and none was made; every reference inside the delivered unit resolves inside it, so no unresolved reference or citation is produced. Printed slips kept literal: no printed word, symbol, hypothesis or number of the counted statement or of its proof was altered. Layout and class: the article's own class is \texttt{amsart} with packages including amsfonts, amsthm, amsmath, amscd, inputenc, t1enc, eucal, indentfirst, graphicx, graphics, pict2e, epic, geometry, epstopdf; the wrapper is the lane's pinned \texttt{amsart} editorial wrapper at 10pt on A4, which restates the theorem-like environments the retained text needs and adds the article's own macro definitions that the retained text needs (none), with extra wrapper packages (bm, bbm, dsfont, xspace, cancel, mathtools). The delivered numbering is the unit's own. Assets: no retained span contains a figure environment, a graphics-inclusion command, a PDF-page inclusion, a table or a TikZ picture; no image file is copied into this package, the package declares no asset dependency and no image file is redistributed with it. Licence: version arXiv:2509.09881v1 of this article is distributed under the Creative Commons Attribution 4.0 International licence, as recorded in the arXiv licence metadata for this exact version and shown on the abstract page https://arxiv.org/abs/2509.09881v1. A full rights notice is delivered at licenses/arxiv-2509.09881-CC-BY-4.0.txt. Characters: every retained excerpt is pure ASCII, so the delivered engine, XeLaTeX with its default Latin Modern text font, reports no missing character.
\begin{Lemma}[Weak extrinsic curvature estimate]\label{lemma_weak extrinsic}
    Let $k>0$ and $H>k$. Given small $\epsilon>0$, there exists a constant $C=C(\epsilon)$ such that the following holds. Let $D$ be an $H$-disk embedded in $\mathbb{H}^3(-k^2)$, and satisfy $0\in D$ and $d_{\mathbb{H}^3(-k^2)}(0,\partial D)\geq\frac{\epsilon}{k}\tanh^{-1}\left(\frac{k}{H}\right)$. Then \begin{equation*}
        |A|(0)\leq CH.
    \end{equation*}
\end{Lemma}

\begin{Cor}[Extrinsic radius estimates]\label{cor_extrinsic radius}
   There exists a constant $R>0$ such that the following estimate holds for any $H>1$. Let $D$ be an $H$-disk embedded in $\mathbb{H}^3$, then we have \begin{equation*}
        \sup_{x\in D}d_{\mathbb{H}^3}(x,\partial D)\leq R\tanh^{-1}\left(\frac{1}{H}\right).
    \end{equation*}
\end{Cor}

\begin{Prop}[Extrinsic curvature estimates]\label{prop_extrinsic curv} 
Given $\delta>0$. There exists a constant $C=C(\delta)$ such that for any $H\geq 1$ and any embedded $H$-disk $D\subset\mathbb{H}^3$,
   we have \begin{equation*}
       \sup_{x\in D: \,d_{\mathbb{H}^3}(x,\partial D)\geq \delta}|A|\leq C.
   \end{equation*}
\end{Prop}

\begin{proof}
First, we use Lemma~\ref{lemma_weak extrinsic} and Corollary~\ref{cor_extrinsic radius} to discuss the case where the mean curvature of $D$ stays away from one and can possibly approach infinity. More specifically, we show the following claim.
\begin{Claim}
    Given $\delta>0$ and $h_0>1$, there exists a constant $C(\delta,h_0)$ such that the following holds. Let $H\geq h_0$ and let $D$ be an $H$-disk embedded in $\mathbb{H}^3$, we have \begin{equation*}
       \sup_{x\in D: \,d_{\mathbb{H}^3}(x,\partial D)\geq \delta}|A|\leq C(\delta,h_0).
   \end{equation*}
\end{Claim}

\begin{proof}[Proof of the claim]
     Arguing by contradiction, the result fails for some $\delta>0$ and $h_0>1$. There exist a sequence of $H_n$-disks $D_n$ in $\mathbb{H}^3$ with $H_n\geq h_0$, and points $x_n\in D_n$ with $d_{\mathbb{H}^3}(x_n,\partial D_n)\geq \delta$, such that \begin{equation}\label{equ_assump_ext_2}
        |A_{D_n}|(x_n)>n.
    \end{equation}
    Rescale the disk by $\Tilde{D}_n:=H_n(x_n)(D_n-x_n)$, then $\Tilde{D}_n$ is a $1$-disk contained in $\mathbb{H}^3\left(-\frac{1}{H_n^2}\right)$. We obtain the following inequalities. \begin{align}\label{equ_1}
         d_{\mathbb{H}^3\left(-\frac{1}{H_n^2}\right)}(0,\partial\Tilde{D}_n) &\geq \delta H_n,\\\nonumber
         |A_{\Tilde{D}_n}|(0)&>\frac{n}{H_n}.
    \end{align}
Let $\epsilon=\frac{\delta}{\tanh^{-1}\left(\frac{1}{h_0}\right)}$. Since the function \begin{equation*}
    x\tanh^{-1}\left(\frac{1}{x}\right)=\frac{x}{2}\ln\left(\frac{x+1}{x-1}\right)
\end{equation*}
decreases on $(1,\infty)$,
the former inequality of \eqref{equ_1} implies that \begin{equation*}
    d_{\mathbb{H}^3\left(-\frac{1}{H_n^2}\right)}(0,\partial\Tilde{D}_n)\geq \delta h_0=\epsilon h_0\tanh^{-1}\left(\frac{1}{h_0}\right)\geq \epsilon H_n\tanh^{-1}\left(\frac{1}{H_n}\right).
\end{equation*}
When combined with Lemma~\ref{lemma_weak extrinsic}, it gives rise to a constant $C'$ that depends only on $\epsilon$, thus depending on $\delta$ and $h_0$, such that 
\begin{equation*}
    |A_{\Tilde{D}_n}|(0)\leq C'.
\end{equation*}
Moreover, Corollary~\ref{cor_extrinsic radius} verifies the existence of $R>0$, so that for the $H_n$-disk $D_n\subset \mathbb{H}^3$, \begin{equation*}
    d_{\mathbb{H}^3}(0,\partial D_n)\leq R \tanh^{-1}\left(\frac{1}{H_n}\right).
\end{equation*}
After rescaling, 
\begin{equation}\label{equ_2}
    d_{\mathbb{H}^3\left(-\frac{1}{H_n^2}\right)}(0,\partial \Tilde{D}_n)\leq RH_n \tanh^{-1}\left(\frac{1}{H_n}\right).
\end{equation}
\eqref{equ_1}, together with \eqref{equ_2}, implies that \begin{equation*}
    \delta\leq R\tanh^{-1}\left(\frac{1}{H_n}\right)\quad \Longrightarrow\quad H_n\leq \frac{1}{\tanh\left(\frac{\delta}{R}\right)}.
\end{equation*}
Consequently, \begin{equation*}
    n\tanh\left(\frac{\delta}{R}\right)\leq \frac{n}{H_n}<|A_{\Tilde{D}_n}|(0)\leq C'.
\end{equation*}
However, since $n\tanh\left(\frac{\delta}{R}\right)$ diverges as $n$ goes to infinity, this contradicts the assumption \eqref{equ_assump_ext_2}. Thus, the proof of the claim is complete.
\end{proof}

To complete the proof of Proposition~\ref{prop_extrinsic curv}, we need to rule out the following scenario. Suppose that $D_n$ is a sequence of $H_n$-disks, where the constants $H_n\geq 1$ converge to one as $n\rightarrow\infty$. Assume there are points $x_n\in D_n$ such that $d_{\mathbb{H}^3}(x_n,\partial D_n)\geq \delta$ and $|A_{D_n}|(x_n)>n$. Consider $\Tilde{D}_n:=|A_{D_n}|(x_n)(D_n-x_n)$. The mean curvature of $\Tilde{D}_n$, given by $\frac{H_n}{|A_{D_n}|(x_n)}$, and the sectional curvature of the ambient space, given by $-\frac{1}{|A_{D_n}|^2(x_n)}$, both converge to zero as $n$ goes to infinity. Consequently, the same argument as in Lemma~\ref{lemma_weak extrinsic} that $\Tilde{D}_n$ converges smoothly on compact subsets to a minimal helicoid in $\mathbb{R}^3$, which leads to a contradiction.
\end{proof}

\end{document}
